\section{Measuring the Score When Growth Is Bought}\label{sec:acquired-growth}
% Generated by code/verify_part2_acquirer.py from notes/case_bending_spoons_f1_extract.md and data/examples/invoice_chain.json (version 2026-10-29c); do not edit.
\providecommand{\acqAcqHOneSix}{\$2{,}286 million}
\providecommand{\acqAcqShareHOneFive}{99.9\%}
\providecommand{\acqAcqShareHOneSix}{99.8\%}
\providecommand{\acqAmzQOne}{11.5\%}
\providecommand{\acqAmzTwentyFive}{11.5\%}
\providecommand{\acqAmzTwentyFour}{12.8\%}
\providecommand{\acqAnnPtsLo}{0.19}
\providecommand{\acqAoimQOne}{51.2\%}
\providecommand{\acqAoimTwentyFiveX}{46.9\%}
\providecommand{\acqAoimTwentyFour}{44.6\%}
\providecommand{\acqCapexHOneSix}{\$4.0 million}
\providecommand{\acqCapexShareHOneFive}{0.05\%}
\providecommand{\acqCapexShareHOneSix}{0.2\%}
\providecommand{\acqCfiQOne}{274.1\%}
\providecommand{\acqCfiTwentyFive}{141.8\%}
\providecommand{\acqCfiTwentyFour}{132.6\%}
\providecommand{\acqChCouponHi}{112.8}
\providecommand{\acqChCouponLo}{86.0}
\providecommand{\acqChPaidHi}{113.1}
\providecommand{\acqChPaidLo}{86.6}
\providecommand{\acqConvHi}{114}
\providecommand{\acqConvLevHi}{103}
\providecommand{\acqConvLevLo}{79}
\providecommand{\acqConvLo}{87}
\providecommand{\acqCouponMarg}{9.625}
\providecommand{\acqDeclHi}{9\%}
\providecommand{\acqDeclLo}{5\%}
\providecommand{\acqDepPrevQOne}{1.5}
\providecommand{\acqDoubleTwentyFive}{8.3}
\providecommand{\acqDoubleZeroTwentyFive}{10.9}
\providecommand{\acqFlagCash}{25\%}
\providecommand{\acqFlagCorr}{17}
\providecommand{\acqFlagCov}{26\%}
\providecommand{\acqFlagCovFast}{42\%}
\providecommand{\acqFlagG}{50\%}
\providecommand{\acqFlagPf}{10\%}
\providecommand{\acqFullPtsHi}{2.2}
\providecommand{\acqFullPtsLo}{0.7}
\providecommand{\acqGiTwentyFive}{\ensuremath{-}5.1\%}
\providecommand{\acqGiTwentyFour}{\ensuremath{-}9.4\%}
\providecommand{\acqGrowthQOne}{132.2\%}
\providecommand{\acqGrowthTwentyFive}{94.7\%}
\providecommand{\acqGrowthTwentyFour}{73.4\%}
\providecommand{\acqIntPaidRatioHOneFive}{90\%}
\providecommand{\acqIntPaidRatioHOneSix}{80\%}
\providecommand{\acqIntPaidShQOne}{14.8\%}
\providecommand{\acqIntPaidShTwentyFive}{10.0\%}
\providecommand{\acqIntPaidShTwentyFour}{2.9\%}
\providecommand{\acqIntPaidShTwentyThree}{4.3\%}
\providecommand{\acqIntQOne}{15.5\%}
\providecommand{\acqIntTwentyFive}{10.9\%}
\providecommand{\acqIntTwentyFour}{4.9\%}
\providecommand{\acqKEb}{1.7}
\providecommand{\acqKVim}{3.2}
\providecommand{\acqKcapHi}{2.1}
\providecommand{\acqKcapLo}{1.2}
\providecommand{\acqLifeHi}{19.5}
\providecommand{\acqLifeLo}{10.6}
\providecommand{\acqLogGrowthTwentyFive}{66.6\%}
\providecommand{\acqMcfUQOne}{24.4\%}
\providecommand{\acqMcfUTwentyFive}{30.5\%}
\providecommand{\acqMcfUTwentyFour}{34.2\%}
\providecommand{\acqNeedTwentyFive}{9\%}
\providecommand{\acqNeedTwentyFour}{11\%}
\providecommand{\acqNotesCoupon}{168}
\providecommand{\acqNrrHi}{95\%}
\providecommand{\acqNrrLo}{91\%}
\providecommand{\acqNrrQOne}{94\%}
\providecommand{\acqNrrTwentyFive}{95\%}
\providecommand{\acqNrrTwentyFour}{91\%}
\providecommand{\acqOcfmQOne}{12.6\%}
\providecommand{\acqOcfmTwentyFive}{22.2\%}
\providecommand{\acqOcfmTwentyFour}{30.5\%}
\providecommand{\acqOpmQOne}{20.0\%}
\providecommand{\acqOpmTwentyFive}{21.3\%}
\providecommand{\acqOpmTwentyFour}{19.0\%}
\providecommand{\acqOutPrevTwentyFive}{2.8}
\providecommand{\acqOutPrevTwentyFour}{2.3}
\providecommand{\acqPfExcess}{100\%}
\providecommand{\acqPfRatio}{2.0}
\providecommand{\acqPfRev}{\$2{,}608.5 million}
\providecommand{\acqPninetyfiveHi}{13}
\providecommand{\acqPninetyfiveLo}{8}
\providecommand{\acqPxVim}{\$1{,}359 million}
\providecommand{\acqQTwoAdjMx}{54.1\%}
\providecommand{\acqQTwoAmz}{11.7\%}
\providecommand{\acqQTwoG}{126\%}
\providecommand{\acqQTwoGx}{126.3\%}
\providecommand{\acqQTwoInt}{15.5\%}
\providecommand{\acqQTwoM}{34\%}
\providecommand{\acqQTwoMx}{34.1\%}
\providecommand{\acqQTwoOrg}{3\%}
\providecommand{\acqQTwoReorgX}{7.2\%}
\providecommand{\acqQTwoRev}{\$704.2 million}
\providecommand{\acqQTwoTrx}{0.9\%}
\providecommand{\acqRAdjQOne}{183.4}
\providecommand{\acqRAdjTwentyFive}{141.6}
\providecommand{\acqRAdjTwentyFour}{118.0}
\providecommand{\acqRAllInvUQOne}{\ensuremath{-}117.5}
\providecommand{\acqRAllInvUTwentyFive}{\ensuremath{-}16.5}
\providecommand{\acqRAllInvUTwentyFour}{\ensuremath{-}25.0}
\providecommand{\acqRGaapQOne}{152.2}
\providecommand{\acqRGaapTwentyFive}{115.9}
\providecommand{\acqRGaapTwentyFour}{92.3}
\providecommand{\acqRIntQOne}{136.7}
\providecommand{\acqRIntTwentyFive}{105.0}
\providecommand{\acqRIntTwentyFour}{87.5}
\providecommand{\acqRLaAvgHiTwentyFive}{111.3}
\providecommand{\acqRLaAvgHiTwentyFour}{97.7}
\providecommand{\acqRLaAvgLoTwentyFive}{104.4}
\providecommand{\acqRLaAvgLoTwentyFour}{94.2}
\providecommand{\acqRLaFullHiTwentyFive}{101.1}
\providecommand{\acqRLaFullHiTwentyFour}{90.5}
\providecommand{\acqRLaFullLoTwentyFive}{92.7}
\providecommand{\acqRLaFullLoTwentyFour}{84.5}
\providecommand{\acqROcfQOne}{144.8}
\providecommand{\acqROcfTwentyFive}{116.9}
\providecommand{\acqROcfTwentyFour}{103.9}
\providecommand{\acqROcfUQOne}{156.6}
\providecommand{\acqROcfUTwentyFive}{125.2}
\providecommand{\acqROcfUTwentyFour}{107.6}
\providecommand{\acqRQTwoAdj}{180.5}
\providecommand{\acqRQTwoGap}{123}
\providecommand{\acqRQTwoOrg}{37}
\providecommand{\acqRQTwoOrgInt}{21.6}
\providecommand{\acqRQTwoOrgIntNoReorgR}{29}
\providecommand{\acqRQTwoOrgIntR}{22}
\providecommand{\acqRQTwoOrgNoReorg}{44.3}
\providecommand{\acqRQTwoOrgNoReorgR}{44}
\providecommand{\acqRQTwoOrgX}{37.1}
\providecommand{\acqRQTwoRep}{160}
\providecommand{\acqRQTwoRepInt}{145.0}
\providecommand{\acqRQTwoRepIntR}{145}
\providecommand{\acqRQTwoRepX}{160.5}
\providecommand{\acqRUcAvgHiTwentyFive}{108.0}
\providecommand{\acqRUcAvgHiTwentyFour}{95.0}
\providecommand{\acqRUcAvgLoTwentyFive}{100.9}
\providecommand{\acqRUcAvgLoTwentyFour}{91.5}
\providecommand{\acqRUcEndHiTwentyFive}{95.2}
\providecommand{\acqRUcEndHiTwentyFour}{86.1}
\providecommand{\acqRUcEndLoTwentyFive}{87.2}
\providecommand{\acqRUcEndLoTwentyFour}{80.3}
\providecommand{\acqReorgQOne}{12.6\%}
\providecommand{\acqReorgTwentyFive}{6.0\%}
\providecommand{\acqReorgTwentyFour}{7.7\%}
\providecommand{\acqRetHi}{28\%}
\providecommand{\acqRetLo}{22\%}
\providecommand{\acqRetNetHi}{22\%}
\providecommand{\acqRetNetLo}{14\%}
\providecommand{\acqRetNetTauHi}{19.8\%}
\providecommand{\acqRetNetTauLo}{13.6\%}
\providecommand{\acqRetTauHi}{25.8\%}
\providecommand{\acqRetTauLo}{22.3\%}
\providecommand{\acqRevQOneSix}{\$601.3 million}
\providecommand{\acqRevTwentyFive}{\$1{,}306.4 million}
\providecommand{\acqRevTwentyFour}{\$671.1 million}
\providecommand{\acqRunRateInt}{572}
\providecommand{\acqRunRateIntHOne}{404}
\providecommand{\acqRunRateIntHi}{787}
\providecommand{\acqRunRateShPf}{14.7\%}
\providecommand{\acqRunRateShPfHi}{20.2\%}
\providecommand{\acqRunRateShQ}{20.3\%}
\providecommand{\acqRunRateShQHi}{27.9\%}
\providecommand{\acqRw}{8.93\%}
\providecommand{\acqRwAnnual}{9.34\%}
\providecommand{\acqTauIres}{24\%}
\providecommand{\acqTrxQOne}{6.3\%}
\providecommand{\acqTrxTwentyFive}{6.5\%}
\providecommand{\acqTrxTwentyFour}{4.8\%}
\providecommand{\acqUcPtsHi}{0.59}
\providecommand{\acqUnpricedDebt}{2{,}234}
\providecommand{\acqVimDays}{327}
\providecommand{\acqVimPre}{\$376.9 million}
% figures generated by code/verify_part2_acquirer.py from the F-1 extract

The frontier of \autoref{prop:frontier} prices growth bought with expensed spend at diminishing returns within one market, and the score budget of \autoref{prop:budget} limits how long a firm can score above its saturated level. A firm that grows by acquiring other firms buys growth with capital instead, across many markets. This section shows how the Rule of $X$ should be measured for such a firm and brackets the score of a serial acquirer that listed in 2026. Section~\ref{sec:data} flags acquisition-affected firm-years and Table~\ref{tab:robust} repeats the main results without them; \autoref{prop:acquired-score} says why such firm-years are not points on the frontier, and \autoref{rem:acquired-bounds} how to score them when they are kept. A companion paper \citep{murthy2026transformation} treats the acquirer's transformation economics; nothing below depends on it. Proofs are in Section~\ref{app:proofs-acquired} of the Supplementary Appendix.

\begin{definition}[Organic and acquired growth]\label{def:acquired-growth}
Revenue grows organically at $g^{\rm org}$, the growth of the businesses the firm already owns, and by acquisition at $g^{\rm acq}$, the revenue run-rate acquired per unit of time as a share of current revenue, so that $g=g^{\rm org}+g^{\rm acq}$; these are instantaneous rates. Measured over year $t$, with $S^{\rm cn}_t$ the revenue of businesses owned throughout the year and $S^{\rm ac}_t$ that of businesses acquired in it, the continuously compounded rates $\ln(S^{\rm cn}_t/S_{t-1})$ and $\ln(1+S^{\rm ac}_t/S^{\rm cn}_t)$ add to $\ln(S_t/S_{t-1})$, and the simple annual rates that the Rule of 40 uses add as $g_t=g^{\rm org}_t+g^{\rm acq}_t$, with $g^{\rm org}_t:=S^{\rm cn}_t/S_{t-1}-1$ and $g^{\rm acq}_t:=S^{\rm ac}_t/S_{t-1}$. Acquisitions cost $k^{\rm acq}$ per dollar of acquired revenue, so the cash paid for them per unit of time, $C^{\rm acq}$, is $k^{\rm acq}g^{\rm acq}$ of current revenue, with $g^{\rm acq}$ the acquired run-rate as a share of current revenue. Over a measured year the charge is $C^{\rm acq}_t/S_t=k^{\rm acq}S^{\rm ac}_t/S_t=k^{\rm acq}g^{\rm acq}_t/(1+g_t)$ when $k^{\rm acq}$ is the price per dollar of revenue recognized in the year, so neither annual measure of acquired growth may be multiplied by $k^{\rm acq}$ directly. The organic score is $R^{\rm org}_X:=Xg^{\rm org}+m$ and the score that charges acquisitions as growth capital is $R^{\rm A}_X:=Xg+m-k^{\rm acq}g^{\rm acq}$. With interest $c_{\rm fin}$ per dollar of revenue, the score after interest is $Xg+m_{op}-c_{\rm fin}$.
\end{definition}

\begin{proposition}[The score when growth is bought]\label{prop:acquired-score}
Under Assumptions~\ref{ass:market}--\ref{ass:capital} and \autoref{def:acquired-growth}:
\begin{enumerate}
\item[(a)] (Net weight.) $R^{\rm A}_X=R^{\rm org}_X+(X-k^{\rm acq})g^{\rm acq}$. Acquired growth raises the score if and only if $k^{\rm acq}<X$; it raises the Rule of 40 if and only if acquired revenue costs less than a dollar per dollar. A score that counts acquired growth without its price, $Xg+m$, overstates $R^{\rm A}_X$ by $k^{\rm acq}g^{\rm acq}$.
\item[(b)] (Bang-bang acquisition.) If revenue can be acquired at the constant price $k^{\rm acq}$ up to the rate $\bar g^{\rm acq}$ and acquisitions leave the growth technology of the businesses already owned unchanged, the maximum of $R^{\rm A}_X$ over spend $\sigma$ and $g^{\rm acq}\in[0,\bar g^{\rm acq}]$ is
\begin{equation}\label{eq:acquired-Rstar}
R^{\rm A*}_X=R^*_X+(X-k^{\rm acq})^+\bar g^{\rm acq},
\end{equation}
with $R^*_X$ of \eqref{eq:Rstar}: the firm buys all the revenue on offer if $X>k^{\rm acq}$ and none if $X<k^{\rm acq}$. Unlike spend, acquisitions meet no diminishing returns until the supply $\bar g^{\rm acq}$ binds.
\item[(c)] (Financing.) The margins of Assumptions~\ref{ass:costs} and \ref{ass:capital} exclude interest. Debt-financed acquisitions leave the operating-margin score unchanged and lower the score after interest by $c_{\rm fin}$. A cash margin from which interest paid has been deducted, as it is from operating cash flow under U.S.\ GAAP, is a margin after interest. A score after interest therefore measures capital structure as well as operations and belongs beside the score before interest and a measure of interest coverage.
\item[(d)] (The existing base.) Let the subscription revenue $S_{\rm sub}$ have net revenue retention $\mathrm{NRR}$, and let the rest of revenue change by the share $\gamma_{\rm ns}$ a year without spend, with shares taken at the start of the year. Before new customers, the businesses owned then contribute $a:=(S_{\rm sub}/S)(\mathrm{NRR}-1)+(1-S_{\rm sub}/S)\gamma_{\rm ns}$ to annual organic growth, so their zero-spend growth, continuously compounded, is $G_i=\ln(1+a)$. It equals $\ln\mathrm{NRR}$ if all revenue is subscription revenue, and lies in $[\ln\mathrm{NRR},0]$ if $\mathrm{NRR}<1$ and $\gamma_{\rm ns}\in[\mathrm{NRR}-1,0]$. By \autoref{rem:assumptions}(b) the saturated score is $R_\infty=m_0+(X-k)G_i$, below $m_0$ when $G_i<0$. New-customer revenue, and so $g^{\rm org}$, is not identified by NRR.
\item[(e)] (Share headroom is bought.) \autoref{prop:budget}, extended by \autoref{rem:assumptions}(b), applies business by business: a business bought at share $s_b$ in its own market can lift its organic score above its saturated level for a total of at most $(X-k)[\ln(\bar s_b/s_b)+(G-G_{i,b})T]$ over $T$ years. If each business sells in a market the firm does not already serve, the firm's budget is the sum over the businesses it owns, so each acquisition buys a fresh budget at $k^{\rm acq}$ per dollar of revenue: the finite headroom of one market becomes a limit on capital and on the supply of targets. An acquisition in a market the firm already serves draws on that market's headroom instead.
\item[(f)] (Treadmill.) Holding growth at $g$ requires $g^{\rm acq}=g+d^{\rm org}$, where $d^{\rm org}:=-g^{\rm org}$ is the organic decay, so $R^{\rm A}_X=Xg+m-k^{\rm acq}(g+d^{\rm org})$, which falls by $k^{\rm acq}$ for each unit of decay. Without new customers, $d^{\rm org}=-G_i$ of part (d), which is $-\ln\mathrm{NRR}$ when all revenue is subscription revenue; if then $G_i<0$, acquisitions must supply more growth than the firm reports.
\end{enumerate}
\end{proposition}

\noindent\emph{Intuition.} Growth bought with capital is the model's growth capital (\autoref{ass:capital}) at the price paid for revenue rather than at $k$. It enters the score at the net weight $X-k^{\rm acq}$, exactly as the net value of a unit of organic growth is $X-k$ (equation \eqref{eq:constants}). A Rule-of-40 screen that adds acquired growth to an operating margin therefore credits every acquisition, whatever it cost; for $X=1$ any purchase above one times its run-rate revenue lowers the instantaneous score that charges its price in full; in an annual score the hurdle also depends on the year's growth and on when the deal closed. Because acquisitions in other markets do not draw down the headroom of the markets the firm is already in, the transience result of \autoref{cor:transient} binds each business, not the firm, and a high score sustained by acquisition is evidence of access to targets and capital, not of a frontier above the ceiling of \autoref{cor:ceiling}. Part (f) is why that access matters: a firm whose base shrinks and that wins no new customers must keep buying to stand still. Acquirer returns are persistent across deals \citep{golubov2015extraordinary}, but bidders that buy often gain on private and subsidiary targets and lose on public ones \citep{fuller2002what}, and a high and uneven acquisition rate lowers performance \citep{laamanen2008performance}, so the supply $\bar g^{\rm acq}$ at a given $k^{\rm acq}$ need not last, and its price may rise with the volume bought.

\begin{remark}[Bounds on the charged score]\label{rem:acquired-bounds}
Let $m^{\rm cf}$ be a cash margin before acquisitions and before financing: operating cash flow plus interest paid net of the tax it saves at the rate $\tau$, over revenue. Under U.S.\ GAAP interest paid is an operating cash flow (ASC 230-10-45-17(d); \citealp{deloitte2026cashflow}), so operating cash flow alone is a margin after interest. A charge on acquisition capital at a cost of capital $r$ that includes the cost of debt, set against it, charges the debt-financed share twice: charge interest or a cost of capital on the same capital, not both. Let $C^{\rm acq}_v$ be the cash paid in year $v$. With the capital recovery factor $\mathrm{CRF}(r,T^{\rm life})=r/(1-e^{-rT^{\rm life}})$, the level annual charge that repays a dollar at rate $r$ over $T^{\rm life}$ years,
\begin{equation}\label{eq:acquired-bounds}
\underbrace{g_t+m^{\rm cf}_t-\frac{C^{\rm acq}_t}{S_t}}_{\text{fully charged}}
\;\le\;
\underbrace{g_t+m^{\rm cf}_t-\mathrm{CRF}(r,T^{\rm life})\!\!\sum_{t-T^{\rm life}<v\le t}\!\!\frac{C^{\rm acq}_v}{S_t}}_{\text{annuity}}
\;\le\;
\underbrace{g_t+m^{\rm cf}_t}_{\text{uncharged}},
\end{equation}
where the first inequality holds whenever the annuity on outlays still within their life is below the current outlay, as it is when outlays grow quickly. A second way to spread an outlay is economic depreciation: if the revenue it buys decays at the rate $d$, the outlay is worth $C^{\rm acq}_ve^{-d(t-v)}$ at $t$, and its user cost is $(r+d)$ times that value \citep{hall1967tax}, as for growth capital that depreciates in \autoref{rem:assumptions}(c). Either charge can be set against capital at the end of year $t$, which charges a full year on outlays made late in it, or against its time-weighted average over the year, which charges each outlay only from the date it is paid and so matches the part of a year's revenue that it adds to $g_t$. Every such charge is nonnegative, so the charged score lies below the uncharged one, and it lies above the fully charged score whenever the year's charge is below the year's outlay. Organic growth is bought with spend that is expensed, so its cost lowers the margin term in the year it is spent ($\sigma$ in \autoref{ass:costs}). Acquired growth is paid for in investing cash flow: it raises the growth term, and on a cash basis its price never touches the score. On the GAAP basis acquisitions reach the score through the amortization and impairment of acquired intangibles, through acquisition-related costs, which are expensed (ASC 805-10-25-23; \citealp{deloitte2026acqcosts}), and through post-combination restructuring; the purchase price itself does not, and goodwill is not amortized \citep{fasb2026asc35020,ifrs2022goodwill}. The uncharged score thus omits the term $k^{\rm acq}g^{\rm acq}$ of part (a). The fully charged score charges each outlay in the year it is paid, as the cash basis charges capital spending $kg$. But a deal closed during year $t$ adds only part of a year's revenue to $g_t$ and the rest to $g_{t+1}$, which is why Section~\ref{sec:data} also flags the following year, so the full price is set against part of the growth it bought: the fully charged score is a lower bound, not an estimate. The annuity and the user cost spread each outlay over the life of the revenue it buys, as depreciation spreads capital spending.
\end{remark}

\noindent\emph{The company's own split.} Bending Spoons S.p.A., which acquires and transforms software businesses and listed on Nasdaq in July 2026 \citep{dealroom2026ipo}, gives the cleanest illustration. In the second quarter of 2026 its revenue grew \acqQTwoG{} on a year earlier, and it put its organic growth at \acqQTwoOrg{} on its own definition, which compares revenue with that of the same businesses a year earlier, using estimated pre-acquisition revenue for businesses bought since \citep{bendingspoons2026sixk}. Unlike \autoref{def:acquired-growth}, that definition counts as organic the growth of an acquired business over its own estimated revenue before the deal. The measure is the company's own: it appears among the unaudited operating results, not among the non-GAAP measures, it is not reconciled, and it rests on estimates. The company does not report a Rule-of-40 score. Computed from its reported revenue growth and its GAAP operating margin of \acqQTwoM, the score is about \acqRQTwoRep, and the organic score $R^{\rm org}_1$ about \acqRQTwoOrg. By part (a), the gap of \acqRQTwoGap{} points, which on the company's definition is the pre-deal revenue of the businesses it bought, is growth credited at the weight $X=1$ with no charge for its price. Two further costs of acquiring sit beside or inside the quarter's margin. Interest on debt raised to pay for acquisitions \citep{bendingspoons2026prospectus}, \acqQTwoInt{} of second-quarter revenue \citep{bendingspoons2026sixk}, lies below the operating margin; after it the organic score is about \acqRQTwoOrgIntR, and the score on reported growth about \acqRQTwoRepIntR. Reorganization expense, \acqQTwoReorgX{} of revenue, which the company attributes to separation packages in the reorganizations of four acquired businesses, lies inside it; treated as a cost of acquiring rather than of the businesses owned, it raises the organic score to about \acqRQTwoOrgNoReorgR{} before interest and \acqRQTwoOrgIntNoReorgR{} after. The organic scores therefore run from about \acqRQTwoOrgIntR{} to \acqRQTwoOrgNoReorgR{} depending on how these two costs are treated. They pair the company's organic growth with a consolidated margin that includes the businesses acquired; they are one quarter, on the company's measure, not our estimate.

\noindent\emph{Reliability of the figures.} The company states in its registration statement and its final prospectus that ``We have identified material weaknesses in our internal control over financial reporting, and we may not be able to remediate these material weaknesses or otherwise maintain effective internal control over financial reporting'' \citep{bendingspoons2026f1,bendingspoons2026prospectus}. No assessment of that control under Section 404 has yet been made, its first-quarter 2026 cash flows are restated, its quarterly reports are unaudited, and its pro forma statements are unaudited and illustrative. We use the figures as filed.

\noindent\emph{Since the second-quarter report.} The scores of Table~\ref{tab:acquired-score} are as of the second-quarter report of 13 August 2026. The company's reports on Form 6-K through 29 October 2026 record the purchases of Airtable and Miro for cash, new and expanded term loans and revolving commitments, new shares sold to Berkshire Hathaway and senior secured notes issued by two subsidiaries to fund future acquisitions \citep{bendingspoons2026sixk}; Section~\ref{app:acquired-after} of the Supplementary Appendix lists them. Interest expense of the first half of 2026 at an annual rate plus the coupons of the new notes is \acqRunRateShQ{} of second-quarter revenue at an annual rate, and \acqRunRateShQHi{} if the new term loans, whose rates the reports do not state, carry the notes' coupon, against the second quarter's \acqQTwoInt. The new debt therefore adds to the interest behind the after-interest scores, which rest on the second quarter. The notes' stated use of proceeds shows continued reliance on acquisition capital, which is why a high charged score is evidence of access to targets and capital; it does not by itself show that the base is shrinking. We take no view on the company's valuation or on any shareholder's investment.

\noindent\emph{Worked numbers.} The company reports its 2024 and 2025 accounts under U.S.\ GAAP in its Form F-1, filed on 8 June 2026 \citep{bendingspoons2026f1}; every F-1 figure we use is unchanged in the final prospectus \citep{bendingspoons2026prospectus}. Table~\ref{tab:acquired-score} recomputes the scores from the full filings, which the replication package contains (Section~\ref{app:data} of the Supplementary Appendix), with simple annual growth rates as the Rule of 40 uses them, and labels each score before or after interest. Revenue grew \acqGrowthTwentyFour{} in 2024 and \acqGrowthTwentyFive{} in 2025 (\acqLogGrowthTwentyFive{} continuously compounded). On the GAAP operating margin, which is before interest, the uncharged scores are \acqRGaapTwentyFour{} and \acqRGaapTwentyFive, and \acqRIntTwentyFour{} and \acqRIntTwentyFive{} after interest by part (c). Operating cash flow is after interest paid, so the cash scores on it, \acqROcfTwentyFour{} and \acqROcfTwentyFive, are scores after interest. Adding back interest net of tax at \acqTauIres, Italy's corporate income tax rate (interest is outside the base of its regional tax; \citealp{pwc2026italy}), gives \acqROcfUTwentyFour{} and \acqROcfUTwentyFive{} before interest. The table adds back interest expense. Adding back interest paid, the cash measure, instead of interest expense moves the 2025 charged scores below to \acqChPaidLo{} to \acqChPaidHi{} (Section~\ref{app:acquired-after} of the Supplementary Appendix). The company's operating cash flows are before deducting stock-based compensation, which the panel's cash score deducts, so its cash scores are not on the panel's basis; its equity compensation expense was \acqSbcTwentyFour{} and \acqSbcTwentyFive{} of revenue in 2024 and 2025, which puts the 2025 cash score after interest paid on the panel's basis, net of equity compensation and capital expenditure, at \acqROcfSbcTwentyFive. On the GAAP basis acquisitions reached the score through amortization and impairment of acquired intangibles, \acqAmzTwentyFour{} and \acqAmzTwentyFive{} of revenue, through transaction expense, \acqTrxTwentyFour{} and \acqTrxTwentyFive, and through reorganization expense, \acqReorgTwentyFour{} and \acqReorgTwentyFive. The company's adjusted operating margin adds all three back, and a score on it would be \acqRAdjTwentyFive{} in 2025: its add-backs are the costs of buying and restructuring businesses, so it lies further from $R^{\rm A}_X$ than the GAAP score. Net investing outflows were \acqCfiTwentyFour{} and \acqCfiTwentyFive{} of revenue. Acquisitions of businesses were \acqAcqShareTwentyFour{} of the net investing outflow in 2024, when purchases of intangible assets (\acqIntangTwentyFour, an asset deal among them, which we count as acquisition capital) were \acqIntangShareTwentyFour, and \acqAcqShareTwentyFiveTwo{} in 2025; in the first halves of 2025 and 2026 they were \acqAcqShareHOneFive{} and \acqAcqShareHOneSix{} of the net investing outflow, and purchases of property and equipment \acqCapexShareHOneFive{} and \acqCapexShareHOneSix{} (\acqCapexHOneSix{} against \acqAcqHOneSix{} in 2026; \citealp{bendingspoons2026sixk}). Charged in full in the year paid, against the margin before interest, the outflows put the score at \acqRAllInvUTwentyFour{} and \acqRAllInvUTwentyFive, a lower bound.

\noindent\emph{Charging acquisitions as capital.} For the charges that spread an outlay over time we take $r=\acqRw$, the continuously compounded value of a \acqRwAnnual{} cost of capital for software firms \citep{damodaran2026cost}, a sector figure below the \acqCouponMarg\% coupon of the company's dollar notes priced on 28 October 2026, and set $d=-\ln\mathrm{NRR}$ and $T^{\rm life}=1/d$, the decay rate and mean life of revenue that decays at the portfolio's retention, which ranged from \acqNrrLo{} to \acqNrrHi{} between 2023 and the first quarter of 2026. This life, \acqLifeLo{} to \acqLifeHi{} years, is a scenario, not an estimate of how long acquired revenue lasts; NRR also covers subscription revenue only. Table~\ref{tab:acquired-score} reports the level annuity and the user cost of \autoref{rem:acquired-bounds}, each against capital at the end of the year and against its time-weighted average, on every net investing outflow since 2023. For the time weights we place 2023 and 2024 outlays at mid-year and 2025 outlays by quarter and half-year from the filings, the second-half outflow between 1 October, the middle of the half, and the close of the year's largest deal, Vimeo, on 24 November 2025 \citep{bendingspoons2026prospectus}. Across the two charges, the two ways of timing them, the range of lives and a tax rate on the interest add-back of zero or \acqTauIres, the 2025 score lies between \acqConvLo{} and \acqConvHi. The same charges set against operating cash flow after interest give \acqConvLevLo{} to \acqConvLevHi{} and charge the debt-financed share twice, by \acqDoubleTwentyFive{} points of 2025 revenue at a \acqTauIres{} tax rate and \acqDoubleZeroTwentyFive{} at zero. The user cost on year-end capital gives the lowest scores, \acqRUcEndLoTwentyFive{} to \acqRUcEndHiTwentyFive{} at a \acqTauIres{} tax rate, because it charges a full year of depreciation and return on capital paid in the last weeks of 2025; time-weighted, the level annuity gives the highest, \acqRLaAvgLoTwentyFive{} to \acqRLaAvgHiTwentyFive. Each point added to $r$ lowers a user-cost score by the ratio of decayed capital to revenue, \acqKcapLo{} to \acqKcapHi{} points in 2025.

The return on the capital is a cleaner statistic than any of these scores, as \autoref{prop:general}(b) suggests. Operating cash flow before interest was \acqRetLo{} to \acqRetHi{} of the time-weighted capital deployed in acquisitions since 2023, written down at the decay rate $d$, across tax rates of zero and \acqTauIres{} on the interest add-back (Table~\ref{tab:acquired-score} uses \acqTauIres), and \acqRetNetLo{} to \acqRetNetHi{} after economic depreciation at $d$, against a cost of capital of \acqRw. Charged at the \acqCouponMarg\% coupon of the company's dollar notes instead of the sector cost of capital, the 2025 charged scores are \acqChCouponLo{} to \acqChCouponHi{} with interest paid added back (\acqChCouponExpLo{} to \acqChCouponExpHi{} with interest expense, the basis of the range above). The ratio counts the cash flow of businesses bought before 2023 but not their price, which overstates it, and leaves the expensed costs of buying and restructuring in the cash flow rather than in the capital, which understates it. Leaving out outlays before 2023, which the extract does not report, makes every charge too small; the company's own measure of capital deployed in acquisitions exceeded the net investing outflow in 2023 and in the first quarter of 2026. Vimeo shows the timing mismatch: it closed on 24 November 2025 for \acqPxVim, so its price was paid in 2025, while \acqVimDays{} days of its 2025 revenue, \acqVimPre, preceded the close and are not in the firm's 2025 revenue. The pro forma statements for 2025, which add AOL, Eventbrite and Vimeo for the full year, show revenue of \acqPfRev, \acqPfRatio{} times actual, and the three purchase prices were \acqKEb{} to \acqKVim{} times annual revenue \citep{bendingspoons2026prospectus}. At those prices $k^{\rm acq}>X=1$, so by part (a) each point of acquired run-rate revenue, as a share of current revenue, lowers the fully charged Rule-of-40 score by \acqFullPtsLo{} to \acqFullPtsHi{} points (in the annual score, which divides growth by last year's revenue, the average full charge per point of acquired growth $g^{\rm acq}_t$ is $k^{\rm acq}/(1+g_t)$, \acqAnnFactor{} times $k^{\rm acq}$ at 2025's growth, so the net effect per point is \acqFullAnnEb, \acqFullAnnAol{} and \acqFullAnnVim{} for the three deals, and lower for deals closed late in the year); charged as a level annuity or at the user cost it costs \acqAnnPtsLo{} to \acqUcPtsHi{} points (\acqAnnPtsAnnLo{} to \acqUcPtsAnnHi{} per point of annual acquired growth) and raises the score. The charging convention, not the data, decides the sign in the year of purchase. For the firm's own decision weight on growth, part (a) gives the sign directly. Charged in full, acquired growth raises $R^{\rm A}_X$ if and only if $X>k^{\rm acq}$. 
The sign flips at $X=k^{\rm acq}$, between \acqKEb{} and \acqKVim{} times revenue for these three deals, so whether a full charge raises or lowers the score depends on the weight the firm acts on, which the data do not identify (Section~\ref{sec:empirics}). Spread as an annuity or a user cost, the charge per point of growth is below one, so acquired growth raises the score at any weight of one or more. The rule that follows from part (a) is to pay less per dollar of acquired revenue than the firm's own weight on growth, which for a value-maximizing firm is, under (H), the marginal value of a dollar of revenue (\autoref{prop:dynamic}).
With NRR of \acqNrrTwentyFive, subscription revenue from customers present a year earlier fell by \acqDeclLo{} in 2025 before new customers, and by \acqDeclHi{} at 2024's \acqNrrTwentyFour{} (\acqGiTwentyFive{} and \acqGiTwentyFour{} continuously compounded). NRR covers subscription revenue only; subscriptions were \acqSubsTwentyFour{} of revenue in 2024 and \acqSubsTwentyFive{} in 2025, but the rate at which the rest of revenue changes without spend is not reported, so the zero-spend growth $G_i$ of part (d) is not identified; it equals $\ln\mathrm{NRR}$ only if all revenue is subscription revenue. Where it is negative, part (d) puts the saturated score below the margin intercept and part (f) says net growth needs new customers or acquisitions. The company put its organic growth for the second quarter of 2026 at \acqQTwoOrg, so on its measure new customers and the rest of revenue more than offset the decline in that quarter. On its own definition the company's organic growth was \acqOrgGTwentyFour{} in 2024 and \acqOrgGTwentyFive{} in 2025 \citep{bendingspoons2026prospectus}; that definition counts the growth of acquired businesses over their estimated pre-deal revenue, so it is not $g^{\rm org}$ of \autoref{def:acquired-growth}, and the split between new customers and acquisitions is not in the filings we read; we do not infer it.

\begin{table}[!tbp]
\centering
\caption{A serial acquirer's scores, recomputed from its F-1 filed on 8 June 2026 \citep{bendingspoons2026f1}, its final prospectus \citep{bendingspoons2026prospectus} and its second-quarter report on Form 6-K \citep{bendingspoons2026sixk} (Rule of 40, simple annual growth; each quarter of 2026 against the same quarter of 2025)}\label{tab:acquired-score}
\footnotesize
\setlength{\tabcolsep}{4pt}\renewcommand{\arraystretch}{0.93}
\begin{tabular}{@{}p{5.7cm}cccc@{}}
\toprule
 & 2024 & 2025 & Q1 2026 & Q2 2026\\
\midrule
Revenue & \acqRevTwentyFour & \acqRevTwentyFive & \acqRevQOneSix & \acqQTwoRev\\
Revenue growth $g$ & \acqGrowthTwentyFour & \acqGrowthTwentyFive & \acqGrowthQOne & \acqQTwoGx\\
Organic growth (company measure, unaudited) & \acqOrgGTwentyFour & \acqOrgGTwentyFive & -- & \acqQTwoOrg\\
GAAP operating margin $m_{op}$ (before interest) & \acqOpmTwentyFour & \acqOpmTwentyFive & \acqOpmQOne & \acqQTwoMx\\
Interest expense $c_{\rm fin}$ & \acqIntTwentyFour & \acqIntTwentyFive & \acqIntQOne & \acqQTwoInt\\
Amortization and impairment of acquired intangibles & \acqAmzTwentyFour & \acqAmzTwentyFive & \acqAmzQOne & \acqQTwoAmz\\
Transaction-related expense & \acqTrxTwentyFour & \acqTrxTwentyFive & \acqTrxQOne & \acqQTwoTrx\\
Reorganization-related expense & \acqReorgTwentyFour & \acqReorgTwentyFive & \acqReorgQOne & \acqQTwoReorgX\\
Adjusted operating margin (company) & \acqAoimTwentyFour & \acqAoimTwentyFiveX & \acqAoimQOne & \acqQTwoAdjMx\\
Operating cash flow / revenue (after interest paid) & \acqOcfmTwentyFour & \acqOcfmTwentyFive & \acqOcfmQOne & --\\
Cash margin before interest $m^{\rm cf}$ & \acqMcfUTwentyFour & \acqMcfUTwentyFive & \acqMcfUQOne & --\\
Net investing outflow / revenue & \acqCfiTwentyFour & \acqCfiTwentyFive & \acqCfiQOne & --\\
Net revenue retention & \acqNrrTwentyFour & \acqNrrTwentyFive & \acqNrrQOne & --\\
\midrule
\multicolumn{5}{@{}l}{\emph{Scores (points)}}\\
Company-adjusted margin, before interest & \acqRAdjTwentyFour & \acqRAdjTwentyFive & \acqRAdjQOne & \acqRQTwoAdj\\
GAAP, before interest: $g+m_{op}$ & \acqRGaapTwentyFour & \acqRGaapTwentyFive & \acqRGaapQOne & \acqRQTwoRepX\\
GAAP, after interest: $g+m_{op}-c_{\rm fin}$ & \acqRIntTwentyFour & \acqRIntTwentyFive & \acqRIntQOne & \acqRQTwoRepInt\\
Organic, before interest: company organic growth $+\,m_{op}$ & \acqROrgTwentyFour & \acqROrgTwentyFive & -- & \acqRQTwoOrgX\\
Organic, after interest & \acqROrgIntTwentyFour & \acqROrgIntTwentyFive & -- & \acqRQTwoOrgInt\\
Organic, before interest and reorganization & \acqROrgNoReorgTwentyFour & \acqROrgNoReorgTwentyFive & -- & \acqRQTwoOrgNoReorg\\
Cash, after interest paid: $g+$ operating cash flow / revenue & \acqROcfTwentyFour & \acqROcfTwentyFive & \acqROcfQOne & --\\
Cash, after interest paid and SBC (panel basis) & \acqROcfSbcTwentyFour & \acqROcfSbcTwentyFive & -- & --\\
Cash, before interest (uncharged): $g+m^{\rm cf}$ & \acqROcfUTwentyFour & \acqROcfUTwentyFive & \acqROcfUQOne & --\\
Charged as capital, before interest (scenario life): & & & & \\
\quad level annuity, year-end capital & \acqRLaFullLoTwentyFour--\acqRLaFullHiTwentyFour & \acqRLaFullLoTwentyFive--\acqRLaFullHiTwentyFive & -- & --\\
\quad level annuity, time-weighted capital & \acqRLaAvgLoTwentyFour--\acqRLaAvgHiTwentyFour & \acqRLaAvgLoTwentyFive--\acqRLaAvgHiTwentyFive & -- & --\\
\quad user cost $(r+d)$, year-end capital & \acqRUcEndLoTwentyFour--\acqRUcEndHiTwentyFour & \acqRUcEndLoTwentyFive--\acqRUcEndHiTwentyFive & -- & --\\
\quad user cost $(r+d)$, time-weighted capital & \acqRUcAvgLoTwentyFour--\acqRUcAvgHiTwentyFour & \acqRUcAvgLoTwentyFive--\acqRUcAvgHiTwentyFive & -- & --\\
Fully charged, before interest (lower bound): $g+m^{\rm cf}-$ outflow / revenue & \acqRAllInvUTwentyFour & \acqRAllInvUTwentyFive & \acqRAllInvUQOne & --\\
\midrule
Return on acquisition capital before / after depreciation & -- & \acqRetTauLo--\acqRetTauHi{} / \acqRetNetTauLo--\acqRetNetTauHi & -- & --\\
\bottomrule
\end{tabular}

\smallskip
\begin{minipage}{0.97\linewidth}\footnotesize
Notes: computed by \texttt{code/verify\_part2\_acquirer.py} from the F-1 and the final prospectus, whose figures agree, and the second-quarter 2026 report and interim statements (unaudited). Basis: GAAP margins are before interest and tax; operating cash flow is after interest paid and cash taxes (U.S.\ GAAP); the cash margin before interest adds back interest expense net of tax at \acqTauIres{} (a zero rate adds a further \acqTauIres{} of $c_{\rm fin}$). Charges: every net investing outflow since 2023, at $r=\acqRw$, $d=-\ln\mathrm{NRR}$ and $T^{\rm life}=1/d$ with NRR from \acqNrrLo{} to \acqNrrHi; the annuity on year-end capital charges each outlay for the full year, the user cost on year-end capital charges $(r+d)$ on capital decayed at $d$, and time-weighted capital charges each outlay from the date it is paid (2023 and 2024 outlays at mid-year; 2025 by quarter and half-year, the second half between 1 October and 24 November). Return: operating cash flow before interest over time-weighted capital written down at $d$, before and after depreciation at $d$ (2025, \acqTauIres{} tax rate). Organic growth is the company's measure, with estimated pre-acquisition revenue. Purchases of property and equipment were \acqCapexRevTwentyFour{} and \acqCapexRevTwentyFive{} of revenue in 2024 and 2025; the panel-basis row deducts them and equity compensation. The Q1 2026 cash flows are as restated in the filing, and the fully charged score for the quarter sets a quarter's outflow against a quarter's revenue. Growth not marked organic includes acquired growth.
\end{minipage}
\end{table}

\noindent\emph{The acquisition flag.} Section~\ref{sec:data} flags a firm-year when cash paid for acquisitions or the increase in goodwill reaches \acqFlagCash{} of prior-year revenue, or pro forma revenue exceeds actual revenue by \acqFlagPf, and it flags the year after as well. The flag covers \acqFlagCov{} of core firm-years and \acqFlagCovFast{} of those growing faster than \acqFlagG. Its cash threshold is a statement about the charge in \eqref{eq:acquired-bounds}: a firm-year flagged on cash paid and growing at $g$ has a fully charged correction of at least $\acqFlagCash/(1+g)$ of revenue, \acqFlagCorr{} points at \acqFlagG{} growth. That is of the order of the \acqPninetyfiveLo{} to \acqPninetyfiveHi{} point fall in the 95th percentile of the score that Table~\ref{tab:robust} reports, so uncharged firm-years of this kind could move the upper tail by as much as the change being measured; the table repeats the results without them and finds the fall within that range. The firm above would carry the flag in every year shown. Its net investing outflows were \acqOutPrevTwentyFour{} and \acqOutPrevTwentyFive{} times prior-year revenue in 2024 and 2025, so it would escape the flag only if less than \acqNeedTwentyFour{} and \acqNeedTwentyFive{} of them paid for acquisitions; in the first half of 2025 acquisitions were \acqAcqShareHOneFive{} of the outflow. The capital it deployed in acquisitions in the first quarter of 2026 alone was \acqDepPrevQOne{} times its 2025 revenue, and its pro forma revenue for 2025 exceeded actual revenue by \acqPfExcess, against the pro forma threshold of \acqFlagPf. Table~\ref{tab:robust} removes such firm-years from the frontier estimates. When they are kept, \eqref{eq:acquired-bounds} is how to score them.

\noindent\emph{What the numbers can and cannot show.} The bounds are wide, from \acqRAllInvUTwentyFive{} to \acqROcfUTwentyFive{} in 2025 on the margin before interest, because the filings split neither growth nor spend by business; charged as capital, the score lies between \acqConvLo{} and \acqConvHi, depending on the charging convention. The final prospectus gives prices and full-year revenue for the three businesses in its pro forma statements, and the company reports its own organic growth for 2024, 2025 and the second quarter of 2026; a split of growth and spend by business is not in the filings we read. What the bounds do show is that the GAAP score of \acqRGaapTwentyFive{} rests on growth whose price the score leaves out, and that its cash counterpart after interest paid, \acqROcfTwentyFive, already carries the interest on debt that paid for part of that growth. A score this high is not evidence that the frontier of \autoref{prop:frontier} can exceed the ceiling of \autoref{cor:ceiling}: by part (e) it is consistent with ordinary frontiers in each acquired business plus capital spent on buying new budgets. The case is one firm, selected because it reached a listing; it has disclosed material weaknesses in its internal control over financial reporting; and its organic growth is measured only by the company, on a definition that counts acquired businesses' growth over their estimated pre-deal revenue.
